\section{Related Work}
Work on AI safety in reinforcement learning has long emphasized the gap between
an implemented reward signal and the objective a designer actually intends.
\citet{amodei2016concrete} catalogued concrete failure modes including reward
hacking and scalable oversight, while AI Safety Gridworlds demonstrated the
value of compact environments for studying such failures under controlled
conditions \citep{leike2017ai}.

Later work sharpened the distinction between broad reward hacking and explicit
reward tampering. Everitt et al.\ analyze tampering incentives with causal
influence diagrams \citep{everitt2021reward}, while Skalse et al.\ characterize
reward hacking as divergence between proxy and true objectives rather than a
grab-bag of anecdotes \citep{skalse2022rewardhacking}. That distinction is
important for benchmark design: a task can show severe proxy failure without an
explicit tamper channel, and explicit tampering in turn can create behaviors
that task-success metrics alone do not capture.

Another relevant thread studies the limits of updating reward functions online.
Even when the reward process is itself adaptive, agents can manipulate weak
oversight or the reward-learning loop \citep{armstrong2020pitfalls}. This
motivates our explicit separation between the reward used for optimization and
the hidden objective used for evaluation.

Procedural generalization work likewise matters here. Layout and nuisance
variation can distinguish brittle reward exploitation from more robust
strategies \citep{cobbe2020leveraging}. In \framework, shift is therefore part
of the task definition rather than an optional stress test. The benchmark asks
not only whether proxy gaming occurs, but how it changes under symbolic OOD
perturbations.

Our contribution is closest to the intersection of these threads: small
model-organism safety environments, formal treatments of reward hacking and
tampering, and benchmark design for robustness. \framework contributes a
compact synthesis centered on explicit task structure, metric discipline, and
leakage-aware monitoring rather than a new optimization technique.
